
\section{Introduction}

Welcome to the manual for the \mcu Extension for \reaper. This
extension enhances support for the Mackie Control Universal protocol.
It was developed specifically for the Mackie Control Universal and the Icon QCon
Pro X, but can also be used for other controllers that use the MCU
protocol. tehsux0r has written a setup guide for the BCF2000 that can
be found \href{http://forum.cockos.com/showthread.php?t=117909}{here}
and nofish wrote a setup guide for the BCR2000 that can be found
\href{http://forum.cockos.com/showthread.php?t=60110}{here}.

\mcu is released under the GPL 3 licence, for details check
section~\ref{disclaimer}.

\subsection{Limitations}
\begin{itemize}
\item \mcu runs on Windows, macOS\,(10.15 Catalina or newer) and
  Linux. The source code is open and is based on the cross-platform
  JUCE library.

\item The minimum supported \reaper version is 6.37. The latest version of the extension for older \reaper versions can be downloaded from \href{https://sync.luckycloud.de/f/591dfbe1cfb042ea8b11/?dl=1}{here}.
\end{itemize}


If these restrictions prevent you from using the extension (and even
if they don't), there are two alternatives:
\href{http://www.mossgrabers.de/Software/Reaper/Reaper.html}{DrivenByMoss
  4 Reaper} and
\href{https://github.com/GeoffAWaddington/CSI_Install}{CSI Project for
  Reaper}.
  

% \subsection{Manual style}
% 
% Buttons are written in \textsc{Small Caps Shape}. Links (to URL-adresses or
% also to different part of the manual) are marked in blue.
% \todo{vervollstaendigen}

\subsection{Most important latest changes}
\subsubsection{Version 0.9.5}
\begin{compactitem}
\item \hyperref[units]{Extender support}: up to eight MCU units (one
  main plus up to seven extenders, i.e.\ up to 64 channels) are
  combined into one logical surface.
\item Cross-platform: in addition to Windows, the extension now also
  runs on macOS and Linux.
\item \hyperref[plugcascade]{FX-Mode page cascade}: with multiple
  units, banks and pages can be spread across the units.
\end{compactitem}

\subsubsection{Version 0.9}
\begin{compactitem}
\item \hyperref[prox]{Icon QCon Pro X support}
\item \hyperref[sendautomation]{Send/Receive automation}
\item Meterbridge improvements
\item New Plug-Mode Options (see Table 4)
\end{compactitem}

\subsubsection{Version 0.8}
\begin{compactitem}
\item 2nd controller support (since v0.8.3; superseded by extender
  support in v0.9.5)
\item \hyperref[buttonactions]{Button-emulating actions} (since v0.8.2)
\item \reaper folders can now be reflected on the MCU using the \foldermode
\item \anchors allows locking a \cs to a fixed track
\item \hyperref[fxpresets]{FX presets} (store and recall FX settings) including
a Blind Test feature
\item \hyperref[localmaps]{FX Local maps} (different maps for each FX
instances of the same FX)
\item improved syncing of the MCU state with the \reaper GUI
\item \hyperref[globalactions]{Global Actions} assignments can be named and
shown in MCU display
\item \hyperref[quickjump]{Quick Jumps} added for faster navigation  
\item improved \hyperref[tracknames]{Track Name} handling
\item FX favorites can be stored in \reaper project file
\end{compactitem}

\subsection{Mailing List / ``Support''}
You can subscribe at \url{https://www.freelists.org/list/mackie_csurf_klinke}
to a mailing list that will be only used to send notifications about
updates.
To communicate with other users you can use the Reaper Forum,
e.g. via this thread
\url{http://forum.cockos.com/showthread.php?t=171827}.
And an issue tracker can be found at \url{https://bitbucket.org/Klinkenstecker/csurf_klinke_mcu/issues?status=new&status=open}.

\subsection{Donations}

Writing the plugin was a lot of work (but also fun most of the time), writing this
manual too (without the fun part). So acknowledgments in the form of donations are
welcome and can be made via PayPal. To donate click
\href{https://www.paypal.com/cgi-bin/webscr?cmd=_s-xclick&hosted_button_id=LR54GZHGL6VHA}{here}
or on the Donate button in the Control Surface
Settings dialog. If your DAW does not have an internet connection, you can also
send the donation via PayPal manually to klinkenstecker@gmx.de. But I am also
delighted when I get positive feedback via email or the \reaper forum ;-)


\subsection{Credits}
I want to thank:
\begin{compactitem}
\item Icon Pro Audio for sending me a QCon Pro X
\item Kathrin for proof-reading the v0.8.0 manual and being such a great
  hiking companion
\item Alex Hayes (plant-on) and Carlos Nunes for commiting bugfixes
  and improving the code
\item Justin for always improving the \reaper API when I need
  additional features (and of course \reaper in general)
\item Chris Doerman (cdoerman) for his Overlay (see
  section~\ref{overlay}), FX maps (e.g. take a look at the hugh Pod
  Farm map, it's an amazing piece of work) and his great help during
  beta-testing
\item Ryan (yagonnawantthatcowbell) and Curvespace for contributing so
  many FX maps
\item axelsk for the Oxford FX maps
\item musicbynumbers for the
  \href{http://forum.cockos.com/showpost.php?p=473466&postcount=127}{Behringer
    BCF2000 setup-guide} and a lot of help on the BCF2000 mode
\item tehsux0r for the even more extensive
  \href{http://forum.cockos.com/showthread.php?t=117909}{Behringer
    BCF2000 setup-guide}
\item nofish for the Behringer
  \href{http://forum.cockos.com/showthread.php?t=60110}{BCR2000
    setup guide}
\item Jason (Drumbum), Miquel (anticlick), Alison and JosMuysers for
beta-testing
\item Jules for the great JUCE library
\item Angeles for all the love
\end{compactitem}

%%% Local Variables: 
%%% mode: latex
%%% TeX-master: "../mcu_klinke_manual"
%%% End: 
